% Created 2018-09-11 Tue 12:03
% Intended LaTeX compiler: pdflatex
\documentclass[11pt]{article}
\usepackage[utf8]{inputenc}
\usepackage[T1]{fontenc}
\usepackage{graphicx}
\usepackage{grffile}
\usepackage{longtable}
\usepackage{wrapfig}
\usepackage{rotating}
\usepackage[normalem]{ulem}
\usepackage{amsmath}
\usepackage{textcomp}
\usepackage{amssymb}
\usepackage{capt-of}
\usepackage{hyperref}
\usepackage[margin=0.80in]{geometry}
\usepackage[backend=biber, style=ieee, url=false]{biblatex}
\usepackage{float}
\usepackage[super,negative]{nth}
\usepackage[capitalise]{cleveref}
\usepackage{pst-node,transparent,ragged2e}
\addbibresource{/home/nuk3/.spacemacs.d/org-files/bibliography.bib}
\DeclareFieldFormat[inproceedings]{citetitle}{\textit{#1}}
\DeclareFieldFormat[inproceedings]{title}{\textit{#1}}
\DeclareFieldFormat[misc]{citetitle}{#1}
\DeclareFieldFormat[misc]{title}{#1}
\renewcommand*{\bibpagespunct}{%
\ifentrytype{inproceedings}
{\addspace}
{\addcomma\space}}
\AtEveryCitekey{\ifuseauthor{}{\clearname{author}}}
\AtEveryBibitem{\ifuseauthor{}{\clearname{author}}}
\SpecialCoor
\author{Nyameko Lisa}
\date{31 August 2018}
\title{IPSP089 - Legal Aspects of Traditional Knowledge and Bioduversity\\\medskip
\large Assignment 1: 781450}
\hypersetup{
 pdfauthor={Nyameko Lisa},
 pdftitle={IPSP089 - Legal Aspects of Traditional Knowledge and Bioduversity},
 pdfkeywords={},
 pdfsubject={},
 pdfcreator={Emacs 27.0.50 (Org mode 9.1.2)},
 pdflang={English}}
\begin{document}

\maketitle
\addvspace{110pt}
\centering{
\pnode(0.5\textwidth,-0.5\textheight){thisCenter}
\rput(thisCenter){%\transparent{0.25}
\includegraphics[width=2.7in]{/home/nuk3/course/llb/wipo-unisa/UNISACoatofArms.eps}}}

\justifying
\addvspace{110pt}
\section*{Declaration}
\label{sec:org9a16b65}
I know that plagiarism is to use someone else’s work and pass it off as my own.\\
I know that plagiarism is wrong.\\
I confirm that this assignment is my own work.\\
I have acknowledged in the bibliography accompanying the assignment all the sources that I have used.\\
I have not directly copied without acknowledgement anything from the Internet or from any other source.\\
I have indicated every quotation and citation in a footnote or bracket linked to that quotation.\\
I have not allowed anyone else to copy my work and to pass it off as their own work.\\
I understand that if any unacknowledged copying whatsoever appears in my assignment I will receive zero per cent for the assignment.\\
I am aware of the UNISA policy on plagiarism and understand that disciplinary proceedings can be instituted against me by UNISA if I contravene this policy.\\
I indicate my understanding and acceptance of this declaration by
entering my name here under:
\begin{itemize}
\item Name: \textbf{Nyameko Lisa} (Student Number: \textbf{7874-909-3})
\end{itemize}

\subsection*{NOTE}
\label{sec:org1d7927d}
Please note that footnotes will be denoted as \footnote{This is a footnote.} and will
appear at the bottom of the page.\\
References will be denoted by \cite{geyer10_towards_clearer_defin_underst_indi} and will appear at the end of the document.
\newpage

\section{Question 1 [50]}
\label{sec:org4937799}

\subsection{Compare and contrast the terms ``traditional knowledge'' and ``indigenous knowledge''. In your answer state whether the \emph{maca} plant falls within these definitions. (10)}
\label{sec:orge906167}
Traditional knowledge systems, creations, innovations, or cultural expressions
refer to \emph{literary, artistic or scientific works, performances, scientific
inventions or discoveries, designs, marks or symbols}, that have evolved in
response to a changing environment and are inherited from one generation to
another within a particular group of people or their territory, \cite{geyer10_towards_clearer_defin_underst_indi}.\\

Although there exists no universally accepted formal definition for indigenous
peoples, `indigenous knowledge' is synonymous with and in fact a subset to,
`traditional knowledge' in terms of its scope, transmission and diversity.
Whilst their distinction is indeed subtle, `traditional knowledge' held by
indigenous people, denotes `indigenous knowledge', \cite{masango10_indigenous_trad_knowl_protect_prospects,amechi15_trad_knowl_medical_uses_plants_pat_sa}.\\

For example, knowledge disseminated through traditional means amongst the
Khoisan peoples of Southern Africa, constitutes `indigenous' or `traditional
knowledge'. In contrast however, the information passed down through traditional
means, by colonial Dutch settlers in the Cape can be described as `traditional
knowledge', but \textbf{not} 'indigenous knowledge'.\\

Given that the maca has long been valued for its ability to enhance fertility in
humans and livestock, it does indeed constitute `traditional knowledge'.
Moreover, given that indigenous peoples have used it the same way for
generations, as a potato-like food source, as well as a supplement for
livestock, appreciated not only for its rich protein and vitamin content, but
for it's sexual performance and fertility enhancement, such use also constitutes
`indigenous knowledge'.

\subsection{How has knowledge been protected customarily? (10)}
\label{sec:org8459a81}

Informal regimes relied on mechanisms based on non-systematic and undocumented
rules, obligations and rights, that are enforced either by elders or religious
leaders. These regimes are `self-legitimised'
\cite{amechi15_leverag_trad_knowl_medicinal}, within and across the community
though the accrual of consensus to abide by and adhere to these laws and
rules.\\

Generally speaking customary law is not recognised by legal institutions outside
of the relevant communities where they were originally established. The culture
and history of a people are preserved by the community elders and leaders, who
are also responsible for determining eligibility and allocation of various arts
to be practised amongst the community's artisans, as well as dispute
resolution.\\

The shortcomings of such protections are clearly evident. Further protection of
ethno-botanical knowledge of the indigenous peoples of Peru will be required,
and in particular will involve the following two aspects:
\begin{itemize}
\item the indigenous peoples of Peru will require protections \emph{`against'} the
intellectual property rights acquired by 'outsider companies' through their
patents for the maca's chemical secrets, which were granted in the United
States, resulting in an appropriation of transitional knowledge, and
\item implementation of the mechanisms of intellectual property legislation to provide
protection of the indigenous peoples of Peru of their traditional knowledge.
\end{itemize}

\subsection{What are the problems confronting rights holders (in our case the people in Junin Peru)? (10)}
\label{sec:orge15fcf6}

Many existing legal frameworks are inadequate at resolving or determining a
common denominator between the economic interests of commercialisation, and the
traditional interests of the people of Junin, Peru. Indigenous knowledge is both
vulnerable to misappropriation under existing industrial property systems and
subsequent financial exploitation by global drug industries,
\cite{masango10_indigenous_trad_knowl_protect_prospects}. All without either
consent from, or acknowledgement of the indigenous peoples, who \emph{arguably} could
be considers the original inventors.\\

Generally speaking, the lack of representation in litigation, coupled with
inadequate protection has resulted in the dismissal and disrespect of the
indigenous people of Peru, who \emph{arguably should} have enjoyed the conferral of
intellectual property rights in relation to the `indigenous knowledge' regarding
the uses of the maca, and again \emph{arguably should} have control over its
exploitation by others, \cite{amechi15_leverag_trad_knowl_medicinal}.\\

Globalisation has given rise to a reluctance of younger generations to learn the
traditional and cultural practices, favouring instead the economic interests of
industrial endeavours. Coupled with the absence of willing heirs and
participants, means that there is significant risk of loss of the traditional
and indigenous knowledge regarding the maca, should it be inadequately
documented and ill-preserved. This problem also manifests itself through the
forced displacement of indigenous traditional communities based on Eurocentric
conservation and biodiversity management strategies, \cite{simelane09_african_trad_knowl_sys_biodiv_manag} \\

Generally speaking modern society holds a prejudiced view, denigrating
traditional knowledge as `primitive'. This lack of appreciation for such
pknowledge, stems from the indigenous people traditionally exploiting the effects
of the maca plant, without describing them in terms of the molecular and
chemical reactions of modern biochemical interactions,
\cite{amechi15_leverag_trad_knowl_medicinal}. This has seen two United States
patents granted for the chemical components, specifically the methods of
extraction form the active ingredients of the maca plant, and not the plant
itself.

\subsection{What is WIPO's role in protecting traditional knowledge? (20)}
\label{sec:org65e4c22}

The WIPO General Assembly established an Intergovernmental Committee on
Intellectual Property and Genetic Resources, Traditional Knowledge and Folklore
focusing primarily on three issues of discussion, and perhaps involving
amendments to \cite{wto17_trips,wipo83_paris_conve_protect_ip} :
\begin{itemize}
\item access to genetic resources and benefit-sharing,
\item the protection of traditional knowledge, innovations and creativity, and
\item the protection of expressions of folklore, including handicrafts.
\end{itemize}

Where the scope of the current work programme of WIPO includes:
\begin{itemize}
\item development of resources detailing methods to implement existing mechanisms of
intellectual property frameworks for the protection of transitional knowledge,
\item complimentary to the above, hands-on national training workshop on existing
intellectual property systems and their use in the protection of traditional
knowledge,
\item with respect to documentation and preservation of transitional knowledge,
establishing mechanisms of their standardisation and recognition as
instruments of intellectual property,
\item fact finding missions, to providing evidence of practical implementations for
the protection of traditional knowledge under existing intellectual property
mechanisms,
\item fact finding missions and feasibility studies into the applicability of
customary laws to traditional knowledge and their integration with existing
statutory regulations,
\item pilot project on the collective acquisition, management, enforcement of
intellectual property rights within traditional knowledge.
\end{itemize}

Unfortunately however, a number of inter-governmental organisational
negotiations with respect to amendments to the TRIPS \cite{wto17_trips} Agreement
on Traditional Knowledge have collapsed,
\cite{saurombe09_protection_indigenous_trad_knowl_through_ip}, and the current
protection afforded is through [article 27(3)(b)]\cite{wto17_trips}, where member
states are empowered to consider protection of traditional knowledge through
intellectual property systems.\\

With respect to protection against exploitation of intellectual property, at the
third session of the Standing Committee on the Law of Patents of WIPO, it was
unsuccessfully proposed that a notification requirement be established. In that
where the subject matter of a patent application is based on the genetic
resources forming part of the ethno-botanical heritage of an indigenous peoples,
then a copy of the contract affording access to the genetic resources in their
country of origin should be filed, requiring possible amendments to \cite{wipo17_pct_ai,wipo70_pct}
.\\

Alternatively, were the indigenous peoples of Peru to have their traditional
knowledge regarding the maca documented and published as identifiable and
searchable prior art, then this would bring into dispute the novelty requirement
of any inventions based on said knowledge. Moreover, this would provide a
mechanism with which the indigenous peoples of Peru could pursue revocations of
the United States patents granted to the companies, albeit at substantial legal
costs. Or perhaps, where such patents are granted, provide co-ownership to
original inventors / indigenous peoples \cite{amechi15_leverag_trad_knowl_medicinal}.\\

Alternatively, rely on the provisions of competition law, where \emph{arguably}
\cite{masango10_indigenous_trad_knowl_protect_prospects}, the maca and any of its
derivative chemical constituent components could be considered a trademark, then
the indigenous peoples of Peru could rely on the protection of trade secrets,
where the patent holders could have said \emph{'trade secrets'} as pertaining to the
maca, disclosed to them in lieu of a license, confidentiality, limited
disclosure, remuneration and royalties, amongst others. Again, said information
will need to be identifiable and searchable, meaning that WIPO will need to
establish mechanisms (such electronic databases\footnote{for example the National
Recordal System (NRS) within the republic, and the Traditional Knowledge Digital
Library (TKDL) in India} \cite{amechi15_leverag_trad_knowl_medicinal}) in order to
strengthen the protection of ethno-botanical knowledge internationally.\\

WIPO's most important role however, together with a number of international,
intergovernmental organisations, will probably be in defining and establishing a
framework for the protection and conservation of indigenous traditional
communities and their biodiversity, \cite{simelane09_african_trad_knowl_sys_biodiv_manag}.
\printbibliography
\end{document}